% \iffalse meta-comment
%
%% regulatory-defs.dtx
%% Copyright 2024-2026 E. Nijenhuis
%
% This work may be distributed and/or modified under the
% conditions of the LaTeX Project Public License, either version 1.3c
% of this license or (at your option) any later version.
% The latest version of this license is in
% http://www.latex-project.org/lppl.txt
% and version 1.3c or later is part of all distributions of LaTeX
% version 2005/12/01 or later.
%
% This work has the LPPL maintenance status ‘maintained’.
%
% The Current Maintainer of this work is E. Nijenhuis.
%
% This work consists of the files listed in the meta-comment of
% regulatory-struct.dtx.
%
% \fi
%
% \iffalse
%<*driver>
\ProvidesFile{regulatory-defs.dtx}
%</driver>
%<package>\NeedsTeXFormat{LaTeX2e}
%<package>\ProvidesPackage{regulatory-defs}
%<package>    [2026/09/10 1.0.0 Xerdi's Regulatory Package (Definitions)]
%
%<*driver>
\documentclass[10pt,english]{ltxdoc}
%! suppress = InclusionLoop
\usepackage{regulatory}
\usepackage{tabularx}
\usepackage[english,dutch]{babel}
%% regulatory-preamble.tex
%% Copyright 2024 E. Nijenhuis
%
% This work may be distributed and/or modified under the
% conditions of the LaTeX Project Public License, either version 1.3c
% of this license or (at your option) any later version.
% The latest version of this license is in
% http://www.latex-project.org/lppl.txt
% and version 1.3c or later is part of all distributions of LaTeX
% version 2005/12/01 or later.
%
% This work has the LPPL maintenance status ‘maintained’.
%
% The Current Maintainer of this work is E. Nijenhuis.
%
% This work consists of the files regulatory.ins,
% regulatory-struct.dtx, regulatory-ref.dtx,
% regulatory-defs.dtx, regulatory-attachments.dtx,
% regulatory-md.dtx,
% regulatory-html.dtx,
% regulatory-english.def, regulatory-dutch.def,
% regulatory.tex, regulatory-preamble.tex,
% regulatory-nl.tex, regulatory-en.tex,
% example1.bib, example2.bib,
% example1-nl.tex, example2-nl.tex,
% example1-en.tex, example2-en.tex,
% md-example.tex, example.md,
% and the derived files regulatory.sty, regulatory-struct.sty,
% regulatory-defs.sty, regulatory-ref.sty,
% regulatory-attachments.sty, regulatory-md.sty
% and regulatory.4ht.
\usepackage{gitinfo-lua}
\usepackage{listings}
\usepackage{adjustbox}
\usepackage{tikz}
\usetikzlibrary{arrows.meta}

\usepackage{fontspec}
\usepackage{textcomp}

\def\projecturl{https://github.com/Xerdi/regulatory}

\definecolor{greycolor}{HTML}{BFBFBF}
\definecolor{primary}{HTML}{174A67}
\definecolor{darkcolor}{HTML}{091D28}
\definecolor{greencolor}{HTML}{176734}
\colorlet{artlabel}{primary}
\colorlet{green}{greencolor}

\lstdefinestyle{tex}{%
    language=[LaTeX]TeX,
    keywordsprefix={\\},
    alsoletter={\\},
    morekeywords={article,para,textfill,masterdocument,refdocument,loadglsdefs,printdefs,setmiddleconjunction,setlastconjunction,setrangeconjunction,setconjunction,markdownInput,glssetwidest,describe}
}

\lstdefinestyle{bash}{%
    language=bash,
    morekeywords={pdflatex,lualatex,latexmk,makeglossaries,bib2gls}
}

\lstdefinelanguage{Markdown}[LaTeX]{TeX}{%
    sensitive=false,
    morecomment=[l][\bfseries\ttfamily\color{green}]{\#},
    morecomment=[s][\bfseries\ttfamily\color{purple}]{:\ \ }{\ },
    morestring=[s][\color{black}]{\{}{\}},
    morekeywords={article,para,textfill,masterdocument,refdocument,loadglsdefs,printdefs,setmiddleconjunction,setlastconjunction,setrangeconjunction,setconjunction,markdownInput,glssetwidest,describe},
    keywordsprefix={\\},
    alsoletter={\\},
}

\lstdefinestyle{md}{%
    language=Markdown
}

\lstdefinelanguage{BibTeX}{%
    keywords={%
        article,book,collectedbook,conference,electronic,entry,ieeetranbstctl,%
        inbook,incollectedbook,incollection,injournal,inproceedings,%
        manual,mastersthesis,misc,patent,periodical,phdthesis,preamble,%
        proceedings,standard,string,techreport,unpublished%
    },
    keywordsprefix={@},
    comment=[l][\itshape]{@comment},
    sensitive=false,
}

\lstdefinestyle{bib}{%
    language=BibTeX,
    morestring=[s][\bfseries\ttfamily\color{purple}]{\ \ \ \ }{\ },
    morecomment=[n][\bfseries\ttfamily\color{green}]{\ \{}{\}},
}

\lstset{
    title=\lstname,
    basicstyle=\ttfamily,
    numberstyle=\ttfamily,
    numbersep=5pt,
    rulecolor=\color{greycolor},
    stringstyle=\color{pink},
    keywordstyle=\color{primary},
    commentstyle=\itshape\color{darkcolor},
    breaklines=true,
    captionpos=b,
    tabsize=4,
    showtabs=true
}

\newcommand\package{\texttt}
\newcommand\file{\texttt}
\newcommand\option{\texttt}

% Reads test/conformance.tex, which test/conformance.sh generates and which is
% kept in the repository: veraPDF is not available everywhere this manual is
% built, so the verdicts travel with the sources.
\newcommand\conformanceversion{?}
\newcommand\conformanceimage{?}
% A digest carries no break points of its own and does not fit a line, so this
% offers one after every character. The argument is expanded first, otherwise the
% loop sees one macro instead of the characters it stands for.
\newcommand\anywherebreakend{}
\def\anywherebreakloop#1{%
    \ifx#1\anywherebreakend\else#1\allowbreak\expandafter\anywherebreakloop\fi}
\newcommand\anywherebreak[1]{\anywherebreakloop#1\anywherebreakend}
\newcommand\conformanceimagetext{%
    \expandafter\anywherebreak\expandafter{\conformanceimage}}

% \ignorespaces, because this macro typesets nothing while its line in the
% generated file still ends in a newline, and that space would otherwise open
% the first cell of the table.
\newcommand\conformancevalidator[2]{%
    \gdef\conformanceversion{#1}\gdef\conformanceimage{#2}\ignorespaces}
\newcommand\conformancerow[6]{%
    \texttt{#1} & \texttt{#2} & \texttt{#3} & \texttt{#4} & \texttt{#5}%
    \ifthenelse{\equal{#6}{}}{}{ \footnotesize(#6)} \\}

% Points at the resulting PDF of an example, which is attached to this manual.
% Its argument is the label of the artifact the example was declared with, and it
% falls back to plain text whenever the PDF isn't there.
\newcommand\exampleresult[1]{%
    \par\noindent
    \translation{The result of this example is}{Het resultaat van dit voorbeeld is}
    \documentattachment{#1}{\file{\artifactfilename{#1}}}.\par}

\def\packagedate{\gitdate}
\def\packageversion{\gitversion}

%\def\cmd{\lstinline[style=tex]}
\def\cmditem#1{\item[\ttfamily\textbackslash{}#1]}

\newcommand\Vtextvisiblespace[1][.3em]{%
  \mbox{\kern.06em\vrule height.3ex}%
  \vbox{\hrule width#1}%
  \hbox{\vrule height.3ex}}

% Package zref-clever loads its language files lazily, which would otherwise
% happen inside a \DocInput region, where the percent sign is ignored and the
% comments of those files end up being parsed as keys. Touching every language
% of this manual schedules them at the start of the document instead, while the
% percent sign still is a comment character.
\zcsetup{lang=english}
\zcsetup{lang=dutch}
\zcsetup{lang=current}

\newcommand\translation[2]{#1}
\begin{document}
    \selectlanguage{english}
    \DocInput{regulatory-defs.dtx}
\end{document}
%</driver>
% \fi
%
% \subsection{\texorpdfstring{\package{regulatory-defs}}{regulatory-defs}}
% \setcounter{CodelineNo}{0}
%
% \subsubsection{\translation{Prerequisites}{Vereisten}}
%
% \package{glossaries-extra} is loaded whichever way the definitions are fed in. Only the \texttt{record}
% option depends on \option{bib2gls}, since that is what hands the recording over to it; everything else
% this module uses -- \cmd{\printunsrtglossary}, \cmd{\glsdefpostlink}, the \texttt{alttree} style -- has
% to be there for the routes that do without \package{bib2gls} as well.
% \iffalse
%<*package>
% \fi
%    \begin{macrocode}
\RequirePackage{regulatory-struct}
\RequirePackage[sanitize=none,nopostdot]{glossaries}
\ifregulatory@bibtogls
\RequirePackage[record,style=alttree]{glossaries-extra}
\else
\RequirePackage[style=alttree]{glossaries-extra}
\fi
%    \end{macrocode}
%
% \subsubsection{\translation{The definitions glossary}{De definitielijst}}
%
% Definitions of the document itself and definitions of other documents are kept apart in two glossary
% types. The \texttt{externals} type is added by \package{regulatory-attachments}.
%    \begin{macrocode}
\newglossary*{definitions}{Definitions}
%    \end{macrocode}
%
% Definition lists are typeset as part of the running text, so the section heading of the glossary is
% dropped altogether. Whenever the list is the first thing in the document, it gets an article heading of
% its own.
%    \begin{macrocode}
\newboolean{regulatory@defslisted}
\setboolean{regulatory@defslisted}{false}
%    \end{macrocode}
%
% Whether the list was printed at all is counted rather than switched, since a counter is global: the list
% may well be printed from inside a group\,---\,a Markdown conversion is one\,---\,and a switch set there
% would be back to false by the end of the document, which is where the question is asked.
%    \begin{macrocode}
\newcounter{regulatory@defsprinted}
\newboolean{regulatory@indexeddefs}
\setboolean{regulatory@indexeddefs}{false}

\renewcommand{\glossarysection}[2][]{\@gobble{#1}\@gobble{#2}}%
%    \end{macrocode}
%
% A listing that opens under a heading opens a line of its own first, and this pulls that line
% back. How much there is to pull back belongs to the style and not to the preamble: the
% \texttt{alttree} that \package{glossaries} brings leaves such a line and the two list styles
% below leave none, so a correction of one line over all three put the first entry of a list on
% the line of the heading itself. A style that needs no correction says so.
%    \begin{macrocode}
\newcommand\regulatory@defsraise{-\baselineskip}
\setglossarypreamble[definitions]{%
    \ifnum\value{article}=0%
        \article{\GetTranslation{Definitions}}\label{art:definitions}%
        \vspace*{\regulatory@defsraise}%
    \else%
        \ifnum\value{parasi}=0%
            \vspace*{\regulatory@defsraise}%
        \fi%
    \fi%
}
%    \end{macrocode}
%
% \subsubsection{\translation{List styles}{Lijststijlen}}
%
% \begin{macro}{regulatory-labeling}
% \begin{macro}{regulatory-description}
% A definition list is a term with its meaning, which is what a description list is for. Two of them are
% offered next to the \texttt{alttree} style \package{glossaries} brings: one built on the
% \texttt{labeling} environment of \package{scrextend}, which aligns every description on the widest name
% given to \cmd{\printdefs}, and one on the \texttt{description} environment of the class, which leaves
% the alignment to the class. Both put the anchor on the name, so a \cmd{\gls} link lands on the term
% rather than beside it.
%
% Which one to reach for is a question of where the document goes. The \texttt{alttree} style is the
% narrowest of the three in HTML: it is written in terms of boxes and widths and comes out as one flat
% paragraph, while both list styles carry the term and its meaning as separate elements.
%    \begin{macrocode}
\newcommand\regulatory@defstyle@body[1]{%
    \renewcommand*{\glossaryheader}{}%
    \renewcommand*{\glsgroupheading}[1]{}%
    \renewcommand*{\glsgroupskip}{}%
    \renewcommand*{\glossentry}[2]{%
        \item[\glstarget{##1}{\glossentryname{##1}}]%
        \glossentrydesc{##1}\glspostdescription}%
    \renewcommand*{\subglossentry}[3]{%
        \glossentrydesc{##2}\glspostdescription}%
}

\newglossarystyle{regulatory-labeling}{%
    \renewcommand\regulatory@defsraise{\z@}%
    \renewenvironment{theglossary}%
        {\begin{labeling}{\@glswidestname}}%
        {\end{labeling}}%
    \regulatory@defstyle@body{}%
}

\newglossarystyle{regulatory-description}{%
    \renewcommand\regulatory@defsraise{\z@}%
    \renewenvironment{theglossary}%
        {\begin{description}}%
        {\end{description}}%
    \regulatory@defstyle@body{}%
}
%    \end{macrocode}
% \end{macro}
% \end{macro}
%
% \begin{macro}{\describe}
% Prints the description of a single definition and marks it as the anchor of that definition, so
% hyperlinks of \cmd{\gls} end up at the right spot.
%    \begin{macrocode}
\newcommand\describe[1]{\setboolean{regulatory@defslisted}{true}\glstarget{#1}{\glsentrydesc{#1}}}
%    \end{macrocode}
% \end{macro}
%
% \begin{macro}{\printdefs}
% Prints the whole definition list, aligned on the width of \meta{width of text}. Within a paragraph the
% list is wrapped in a \texttt{minipage}, to keep it from being broken across pages halfway an enumeration.
%
% The width is read before it is set, since \cmd{\glssetwidest} stores it in the very macro one would
% reach for to name the widest name already known, and setting that macro to itself is a loop the run does
% not come back from.
%
% The optional argument picks the style. A name for which a \texttt{regulatory-} style exists is taken
% to mean that one, so \texttt{labeling} and \texttt{description} reach the two above; anything else is
% handed to \package{glossaries} as it stands, which leaves every style it knows available. Without the
% argument the \option{defstyle} option decides.
%
% Which of the two printing commands applies depends on how the entries arrived rather than on the
% \option{bib2gls} option. \cmd{\printglossary} needs a sorted and indexed glossary, which only
% \cmd{\loadglsentries} produces; \package{bib2gls} and \cmd{\newdefinition} both hand over entries
% that are already in the order they should be printed in, and \cmd{\printunsrtglossary} prints those
% without an external tool having to run at all.
%    \begin{macrocode}
\newcommand\printdefs[2][\regulatory@defstyle]{%
    \renewcommand*{\glsnamefont}[1]{\textmd{##1}}%
    \setboolean{regulatory@defslisted}{true}%
    \stepcounter{regulatory@defsprinted}%
    \protected@edef\regulatory@thiswidest{#2}%
    \expandafter\glssetwidest\expandafter{\regulatory@thiswidest}%
    \edef\regulatory@thisstyle{%
        \ifcsname @glsstyle@regulatory-#1\endcsname regulatory-#1\else#1\fi}%
    \ifnum\value{parasi}>0%
        \def\@@wrapper##1{\begin{minipage}{\linewidth}##1\end{minipage}}%
    \else%
        \def\@@wrapper##1{##1}%
    \fi%
    \@@wrapper{%
        \begingroup%
        \setglossarystyle{\regulatory@thisstyle}%
        \ifthenelse{\boolean{regulatory@indexeddefs}}{%
            \printglossary[type=definitions,title={},nonumberlist=true]%
        }{%
            \printunsrtglossary[type=definitions,title={},nonumberlist=true]%
        }%
        \endgroup%
    }%
}
%    \end{macrocode}
% \end{macro}
%
% \begin{macro}{\loadglsdefs}
% Loads the definitions of \meta{src} under the \texttt{definitions} type. With \package{bib2gls} the
% \option{alldefs} option is the difference between listing every entry of the resource and only the used
% ones.
%    \begin{macrocode}
\edef\@regulatory@bib@selection{}
\ifregulatory@alldefs
\edef\@regulatory@bib@selection{,selection=all}
\fi

\newcommand\loadglsdefs[1]{%
    \setboolean{regulatory@defslisted}{true}%
    \ifregulatory@bibtogls%
        \GlsXtrLoadResources[type=definitions,src={#1},sort={nl-NL},category={defs}\@regulatory@bib@selection]%
    \else%
        \setboolean{regulatory@indexeddefs}{true}%
        \loadglsentries[definitions]{#1}%
    \fi%
}
%    \end{macrocode}
% \end{macro}
%
% \begin{macro}{\newdefinition}
% Declares one definition in the document itself, under the \texttt{definitions} type and the
% \texttt{defs} category, so that it lands in the same list as the ones read from a file. This is the
% route that needs no second file and no second program: the entries are printed in the order they are
% declared in.
%    \begin{macrocode}
\newcommand\newdefinition[3]{%
    \setboolean{regulatory@defslisted}{true}%
    \newglossaryentry{#1}{type=definitions,category=defs,name={#2},description={#3}}%
}
%    \end{macrocode}
% \end{macro}
%
% Only the indexed route asks for a glossary to be made, which is what draws in \texttt{makeindex} or
% \texttt{makeglossaries}; declaring the entries in the document or reading them with \package{bib2gls}
% does not. With the \option{alldefs} option every unused definition is added at the end of the document,
% so that it shows up in the list as well.
%
% The kernel hook rather than \cmd{\AtEndPreamble} of \package{etoolbox}, which is what this was and which
% cost a \package{xpatch} in the prerequisites for one call. Both fire in the same place; the hook is the
% one that is there without asking for it.
%    \begin{macrocode}
\AddToHook{begindocument/before}{%
\ifthenelse{\boolean{regulatory@indexeddefs}}{\makeglossaries}{}%
}
\AtEndDocument{%
\ifregulatory@alldefs\glsaddallunused\fi%
}
%    \end{macrocode}
% \iffalse
%</package>
% \fi
%
% \Finale
%
